\section{A canonical detected quotient of class three}
\label{sec:canonical-class-three}

We enlarge the canonical quotient by its central relations.
Let \(q_2=q:H\twoheadrightarrow Q_2=Q\) be the quotient of
Theorem~\ref{thm:canonical-torsion-quotient}. Write
\[
 Q_2=(W\oplus\overline{\mathcal E},\star),\qquad
 \overline{\mathcal E}=\mathcal E/Z,\qquad
 N=\widehat{\bigwedge}^{\,2}_{\mathbf Z_p}W\subseteq\overline{\mathcal E}.
\]
Thus \(W=U/D\), \(D=U_{\Lambda\text{-tors}}\), and
\[
 0\longrightarrow N\longrightarrow\overline{\mathcal E}
 \longrightarrow T/\pi(Z)\longrightarrow0.
\]
Define, using closed commutator subgroups,
\[
 B=\ker(H\to W),\qquad K_2=\ker q_2,\qquad
 J_4=\overline{[H,K_2][B,B]},\qquad
 \mathcal R=K_2/J_4.
\]
Let \(K_3\) be the inverse image of \(\mathcal R_{\Lambda\text{-tor}}\) in \(K_2\), and put \(L^{(2)}=\Omega_p^{K_2}\).

Since \(q_2(B)=\overline{\mathcal E}\) is central,
\([H,B]\subseteq K_2\), and \(\mathcal R\) is central in \(H/J_4\).
The full arithmetic group preserves all these subgroups.
Its action on \(\mathcal R\) factors through \(\Gamma\times\Delta\).

\begin{theorem}[Canonical class-three detection]
\label{thm:canonical-class-three}
For every odd prime $p$, the module
$\mathcal R_{\Lambda\text{-tor}}$ is closed and $K_3$ is
normal in $\operatorname{Gal}(\Omega_p/\mathbb Q)$.
The pro-$p$ outer-action kernel in $H$ is contained in $K_3$.
Consequently
\[
 L^{(2)}\subsetneq L^{(3)}=\Omega_p^{K_3}\subseteq\Omega_p^*,
 \qquad \operatorname{Gal}(L^{(3)}/F_\infty)=Q_3=H/K_3.
\]
The group $Q_3$ is torsionfree and has nilpotency class exactly three.
There is a natural injection
\[
 \operatorname{Lie}^3(W)\lhook\joinrel\longrightarrow
                         \ker(Q_3\to Q_2).
\]
In particular, every nonzero degree-three Lie word in the odd
quotient $W$ survives.
\end{theorem}

\begin{proof}[Proof of Theorem~\ref{thm:canonical-class-three}]
Put \(\mathcal S=[H,B]/J_4\). We first calculate this central module
integrally:
\begin{equation}\label{eq:class3-S}
 \mathcal S\simeq
 \frac{W\widehat\otimes_{\mathbf Z_p}\overline{\mathcal E}}
 {J(\widehat{\bigwedge}^{\,3}W)},\qquad
 J(x\wedge y\wedge z)
 =x\otimes(y\wedge z)+y\otimes(z\wedge x)+z\otimes(x\wedge y).
\end{equation}
The map is the commutator of lifts of \(W=H/B\) and
\(\overline{\mathcal E}=B/K_2\). Changing either lift changes it by
\([B,B]\) or \([H,K_2]\), so it is well-defined.
It is bilinear because its image is central, and Hall--Witt
gives the displayed Jacobi relation.

For the exact presentation, split \(U=D\oplus W\) as compact
\(\mathbf Z_p\)-modules. Choose a convergent basis of \(W\),
finite bases of \(D,T\), and lifts forming a free pro-\(p\)
basis of \(H\). The lifts may be chosen so that \(D\) maps
to \(1\), the \(W\)-basis maps to the odd section of \(Q_2\),
and the \(T\)-basis maps to selected central lifts.
The necessary corrections follow from the surjection
of the derived subgroup onto \(N\).

Split \(\mathcal E=N\oplus\bigoplus_i\mathbf Z_pt_i\), and write a
basis of \(Z\) as
\[
 z_s=n_s+\sum_i a_{si}t_i.
\]
Modulo \(J_4\), the \(D\)-lifts are central,
\([W,t_i]\) and triple \(W\)-commutators are central,
and \([t_i,t_j]\), \([t_i,[W,W]]\), and brackets involving
\(D\) vanish. The surviving central derived generators
have exactly the relations
\[
 \text{Jacobi},\qquad
 [W,z_s]=[W,n_s]+\sum_i a_{si}[W,t_i]=0.
\]
Indeed, collect a free nilpotent group of class three in a
Hall basis. Before the \(z_s\)-relations, its surviving
central module is
\[
 \operatorname{Lie}^3(W)
 \oplus\left(W\widehat\otimes\bigoplus_i\mathbf Z_pt_i\right).
\]
The remaining relations are precisely \(W\widehat\otimes Z\),
which proves \eqref{eq:class3-S}. No coefficient \(a_{si}\)
has been inverted. In particular all integral torsion
arising from \(T/\pi(Z)\) remains in this presentation.
The Hall collection and the leading Magnus terms give
the integral independence of these generators; passing
through finite-coordinate quotients gives the compact
statement. This uses no class-three BCH denominator and
is valid also at \(p=3\).

The abelian image gives the further exact sequence
\begin{equation}\label{eq:class3-R}
 0\longrightarrow \mathcal S\longrightarrow \mathcal R
 \longrightarrow D\oplus\pi(Z)\longrightarrow0.
\end{equation}
To identify its kernel, the derived subgroup of
\(H/[H,B]\) is \(\widehat{\bigwedge}^2W=N\), and the induced
map to \(Q_2\) is injective there. Hence
\(K_2\cap[H,H]=[H,B]\).
The image of \(K_2\) in \(H^{\rm ab}\) is exactly the
kernel of \(M\to Q_2^{\rm ab}\), namely \(D\oplus\pi(Z)\).
This proves \eqref{eq:class3-R}.

Let \(p^b\) kill the finite torsion subgroup of \(T/\pi(Z)\).
We claim
\begin{equation}\label{eq:class3-pbound}
                         p^b\mathcal S[p^\infty]=0.
\end{equation}
Let \(N^{\rm sat}\) be the \(p\)-saturation of \(N\) in
\(\overline{\mathcal E}\). Then \(p^bN^{\rm sat}\subseteq N\), and
\(\overline{\mathcal E}/N^{\rm sat}\) is torsionfree.
The Jacobi image is saturated in \(W\widehat\otimes N\),
because its quotient is the torsionfree module
\(\operatorname{Lie}^3(W)\). Explicitly, in each component
with three distinct basis indices its single relation is
primitive, with coefficients \(1,1,1\) up to sign; repeated
indices give the usual free components.
If \(p^kx\in J\) for \(x\in W\widehat\otimes\overline{\mathcal E}\), its
image in \(W\widehat\otimes(\overline{\mathcal E}/N^{\rm sat})\) vanishes.
Thus \(p^bx\in W\widehat\otimes N\), and saturation there
gives \(p^bx\in J\). This proves \eqref{eq:class3-pbound}.
This bounds the exponent of the entire compact $p$-torsion subgroup.

We next prove that the target length-three map is
injective on \(\mathcal S[1/p]\).
No integral \(\Lambda\)-basis of \(W\) is required.
Choose the established \(\Lambda[1/p]\)-basis of
\(W[1/p]\); finite generation gives a bounded lattice
\(p^a\Lambda^r\subseteq W\subseteq p^{-a}\Lambda^r\).
Every fixed completed tensor degree inherits a common
denominator bound. The integral presentation above used
only a compact \(\mathbf Z_p\)-basis.
Use the cofinite free rational target already chosen for
\(q_2\). The image of \(B\) starts in generator length two,
and the image of \(K_2\) starts in length three.
Therefore \([B,B]\) and \([H,K_2]\) vanish at lengths at
most three, giving a map from \(H/J_4\).
Its restriction to \(\mathcal R\) is additive and has pure length
three.

Over a finite splitting field \(L/\mathbf Q_p\), choose
Jordan blocks for \(\gamma\) on
\(T'=(T/\pi(Z))[1/p]\).
For a block of eigenvalue \(\alpha_i\) and length \(\ell_i\),
choose a lift \(y_i\in\overline{\mathcal E}_L\) of its top vector and put
\[
 C_i=(\gamma-\alpha_i)^{\ell_i}y_i\in N_L.
\]
For each fixed \(\alpha\), the residues of the \(C_i\) with
\(\alpha_i=\alpha\) are linearly independent over \(L\) in
\(N_L/(\gamma-\alpha)N_L\). Otherwise a relation
\(\sum_i d_iC_i=(\gamma-\alpha)A\) would make
\[
 \sum_i d_i(\gamma-\alpha)^{\ell_i-1}y_i-A
\]
a \(\gamma\)-eigenvector in the \(\Lambda\)-torsionfree module
\(\overline{\mathcal E}_L\). Its projection is a linear combination of
independent Jordan socle vectors, so all \(d_i\) must vanish.

In bounded pair coordinates the target kernel of
\((\gamma-\alpha_i)^jy_i\), \(0\le j<\ell_i\), is
\[
 \frac{C_i(X,Y)}
 {((1+X)(1+Y)-\alpha_i)^{\ell_i-j}}.
\]
All pair evaluation denominators were excluded in the
cofinite target chosen for \(Q_2\).
The source \eqref{eq:class3-S}, rationally, is the direct sum
of \(\operatorname{Lie}^3(W)_L\) and the completed tensor
of \(W_L\) with these finitely many lifted Jordan vectors.

Suppose an element has zero target length-three image.
The moment identity, after clearing denominators, gives
an identity in three independent position variables.
Choose a largest nonzero mixed pole of order \(k\) at
the pair \((1,2)\), with root \(\alpha\), and use
\[
 V=(1+T_1)(1+T_2)-\alpha.
\]
Poles at \((2,3)\) are units in the \(V\)-series field:
their constant terms are
\(\alpha(1+T_3)/(1+T_1)-\alpha'\ne0\).
Different roots at \((1,2)\) are units as well.
The coefficient of \(V^{-k}\) is consequently
\[
 \sum_i F_i(T_3)
 C_i\!\left(T_1,\frac{\alpha}{1+T_1}-1\right)=0
\]
as an array over pair colors, for each singleton color.

Constant-field independence suffices here. Choose finitely
many coefficient-extraction functionals in \(T_1\) and
the pair colors giving an invertible matrix over \(L\)
on the finite list of residue vectors.
Apply them coefficientwise in the disjoint variable
\(T_3\). The resulting invertible constant matrix forces
all \(F_i(T_3)=0\), contradicting the chosen pole.
No independence over the pair-variable function field
has been assumed. Thus all mixed terms vanish.
The remaining pure odd triple is detected by the completed
triple moment map.

This proves injection of \(\mathcal S[1/p]\) into rational
independent-position coefficient arrays before evaluation.
On those arrays, a nonzero \(f\in\Lambda\) acts by the
nonzero multiplier
\[
 f((1+T_1)(1+T_2)(1+T_3)-1).
\]
Setting two variables to zero proves nonvanishing.
Hence \(\mathcal S[1/p]\) is \(\Lambda\)-torsionfree.

Choose a characteristic polynomial \(a(z)\) annihilating
\(\gamma\) on the finite free module \(D\oplus\pi(Z)\).
By \eqref{eq:class3-R}, \(a(\gamma)\mathcal R\subseteq \mathcal S\).
If \(r\in \mathcal R_{\rm tor}\), then \(a(\gamma)r\) is torsion
in \(\mathcal S\), so its rational image is zero. Equation
\eqref{eq:class3-pbound} gives \(p^ba(\gamma)r=0\).
Conversely this nonzero \(\Lambda\)-annihilator makes
\(r\) torsion. We have proved the explicit identity
\[
 \mathcal R_{\rm tor}
 =\ker\bigl(p^ba(\gamma):\mathcal R\longrightarrow \mathcal R\bigr).
\]
Thus \(\mathcal R_{\rm tor}\) is closed and arithmetic-stable.
The quotient \(\mathcal R/\mathcal R_{\rm tor}\) is torsionfree, and the
extension
\[
 1\longrightarrow \mathcal R/\mathcal R_{\rm tor}
 \longrightarrow Q_3\longrightarrow Q_2\longrightarrow1
\]
makes \(Q_3\) torsionfree as well.

The pro-$p$ outer kernel lies in \(K_2\).
For its image \(r\in \mathcal R\), the target length-three image
of \(a(\gamma)r\in \mathcal S\) is zero. The rational injection
shows this element is \(p\)-torsion, so
\(p^ba(\gamma)r=0\). Therefore \(r\in \mathcal R_{\rm tor}\), proving
outer-kernel containment in \(K_3\) and the claimed field
inclusions.

Finally the inclusion \(N\hookrightarrow\overline{\mathcal E}\) and
\eqref{eq:class3-S} give
\[
 \operatorname{Lie}^3_{\mathbf Z_p}(W)\hookrightarrow \mathcal S.
\]
Its triple coefficient module is \(\Lambda\)-torsionfree,
so it meets \(\mathcal R_{\rm tor}\) trivially and survives in
\(\ker(Q_3\to Q_2)\).
It is nonzero, since \(W\) has infinite \(\mathbf Z_p\)-rank.
Thus \(Q_3\) has class exactly three and the extension over
\(Q_2\) is strict.
\end{proof}

